\chapter{Inside a Humanvoice review}
\label{ch:product}

The previous chapter followed a research note from its muddled opening to a
reader's decision. We now return to the same teaching case and move closer to
the screen. One paragraph is enough to show what the author, technical reviewer,
and reader would actually receive. Its text and numbers remain synthetic: they
expose the product's duties without masquerading as evidence from a live trial.

The same version appears differently to the author, technical reviewer, and
reader. Figure~\ref{fig:three-views} shows how one record supports those views
without exposing the private review trail in the manuscript.

\hvFigureThreeViews

\section{Before a reader sees a revision}

The expensive mistake is to give a senior reader a document that was never
ready for a senior reader. Humanvoice therefore begins before any prose changes.
The author states who will read the revision, what the reader already knows,
which objects may not change silently, and what voice the genre requires, then
invokes \texttt{hv humanize} on a finished manuscript. That invocation is the
authorization to change text. The manuscript itself is immutable: the system
copies it and writes a separate child revision.

The first output is a concept inventory over the source, not a plan for a new
document. Humanvoice partitions every reader-facing span, records each
substantive concept it carries---definitions, distinctions, claims, mechanisms,
assumptions, derivation steps, qualifications, examples, counterexamples,
evidence interpretations, and teaching transitions---and states what each
concept must do for this reader. A second, independent pass reads the inventory
back against the source to catch anything the first pass missed. Ambiguities go
to a human, who adjudicates them and freezes the baseline. Everything downstream
is measured against that frozen baseline.

Consider the teaching case in this chapter. The source paragraph carries several
concepts: the comparison between the two models, the funding wedge and its unit,
which likelihood terms and observations are compared, what agreement would and
would not authorize, and the limit of the calibration. Humanvoice records each
one, notes that the wedge is stated without its unit and that the calibration
limit is asserted without a mechanism, and assigns each concept the explanation
it owes the reader:

\begin{enumerate}
\item define the two models before comparing them;
\item state the funding wedge with its unit and location;
\item name which likelihood terms, observations, and paths are compared;
\item derive why agreement authorizes the next step, and where it stops; and
\item interpret the diagram against the objects it connects.
\end{enumerate}

Only after the author accepts that inventory does the writer rewrite one source
passage at a time. Each unit receives the exact source span, its assigned
concepts and obligations, the protected objects it contains, and enough
neighboring context to keep transitions coherent. A unit that cannot discharge
its obligations within one call is split on concept boundaries and resumed. It is
never trimmed to fit. This is slower for the machine and faster for the person
who would otherwise have to reconstruct the argument while reading.

\begin{figure}[H]
\centering
\resizebox{0.98\textwidth}{!}{%
\begin{tikzpicture}[font=\small,node distance=5mm and 4mm]
  \node[hvstage,text width=25mm] (source) {Immutable source\\\footnotesize finished manuscript};
  \node[hvstage,right=of source,text width=25mm] (plan) {Frozen concept baseline\\\footnotesize concepts, obligations, dependencies};
  \node[hvgoldbox,right=of plan,text width=25mm] (draft) {Source-grounded rewrite\\\footnotesize one passage at a time};
  \node[hvwarn,right=of draft,text width=25mm] (check) {Semantic preflight\\\footnotesize retention, explanation, integrity};
  \node[hvbluebox,right=of check,text width=25mm] (packet) {Reader packet\\\footnotesize only after every gate passes};
  \node[hvbluebox,right=of packet,text width=25mm] (reader) {Human reader\\\footnotesize one expensive final decision};
  \draw[hvflow] (source)--(plan)--(draft)--(check)--(packet)--(reader);
  \draw[hvback] (check.south) to[bend right=28] node[below,hvnote]{failed gate: located repair and rebuild} (draft.south);
  \node[hvguard,fit=(check)(packet),inner sep=3mm] {};
\end{tikzpicture}}
\caption{Humanvoice spends scarce reader attention at the end of the path. A
failed check returns the revision to a located repair; it cannot produce a
reader packet by default.}
\label{fig:pre-human-pipeline}
\end{figure}

The guard in Figure~\ref{fig:pre-human-pipeline} is a product contract, not a
promise that a machine can certify taste. It says something narrower and more
useful: an unfrozen baseline, a source concept without verified correspondence in
the revision, a concept named but not explained, an unsupported addition, a
concept used before its prerequisite, an unexplained protected-object change, a
broken build, or an unresolved finding cannot be packaged as ready for a human
reader. The system may still keep an internal revision for repair. It does not
ask an expensive reader to serve as its debugger.

\section{A draft that should never reach the reader}

The note compares the baseline model with the constrained reconstruction. An
unguarded writer has produced the paragraph below. In the old workflow it
would have been sent to the owner for diagnosis. In the proposed workflow it is
a fixture: the pre-human checks should locate the problem before a reader sees
the document.

\begin{quote}\small
"In the Phase-2 rescue, the declared restricted reconstruction is translated
into the solver and compared like with like. We do not claim that this proves
the policy works in general; rather, it is a robust and transparent framework
for checking the relevant objects, including the 35 basis-point funding wedge
in Table~3 and the likelihood terms in Equation~(12). The result is therefore
useful for stakeholders who need confidence that the pipeline is working."
\end{quote}

The paragraph is not nonsense. It has a real comparison, a numerical result,
and an appropriate desire not to overstate. It is nevertheless difficult for a
reader arriving without the project history. "Phase-2 rescue" names a private
event. "Like with like" does not identify the objects being compared. The
negative definition delays the positive mechanism. "Stakeholders" is an
abstract actor. The reader cannot tell whether the 35 basis points are an
estimated effect, a calibration input, or an illustrative value.

The paragraph fails several baseline and preflight checks. The private phase
name is absent from the approved vocabulary; the comparison has no named
objects; the number has no stated role; and the final sentence names no action.
Those are located failures, not invitations to score the whole document. The
writer receives the evidence and a repair question while the source remains
unchanged. If the question cannot be answered from the brief and evidence
ledger, the run stops rather than inventing a smoother claim.

\section{The review record}

Before anyone edits the paragraph, the parser maps its words to source spans
and gives the equation reference, table reference, number, and declared result
protected identities. The first review then records five findings. The table on
the next page makes each one contestable.

\begin{table}[H]
\centering
\small
\begin{tabularx}{\textwidth}{@{}p{0.09\textwidth}p{0.22\textwidth}L p{0.22\textwidth}@{}}
\toprule
ID & Location & Finding and evidence & Editorial question \\
\midrule
F1 & Sentence 1 & "Phase-2 rescue" is a private status label; no source
definition appears in the document. & What event does the reader need to know,
and can it be named directly? \\
F2 & Sentence 1 & "Like with like" names a comparison without naming its
objects. & Which likelihood terms, observations, or paths are compared? \\
F3 & Sentence 2 & The negative definition occupies the first clause and the
positive mechanism is delayed. & What does the analysis establish in this
document? \\
F4 & Sentence 2 & The number 35 and the references to Table~3 and Equation~(12)
are protected. & Are the unit, sign, and interpretation of the wedge explicit? \\
F5 & Sentence 3 & "Useful for stakeholders" asserts value without naming a
decision or evidence. & Which reader action would change if the comparison
passes? \\
\bottomrule
\end{tabularx}
\caption{A finding is located, reasoned, and contestable. It is not an automatic
rewrite instruction.}
\label{tab:worked-findings}
\end{table}

F1, F2, and the source-location part of F4 are deterministic. A parser can
prove that a term is undefined in the current document. F3 and F5 remain
editorial questions because the author must decide whether a qualification is
necessary for an honest conclusion.

The author can reject F1 if the private label is deliberately part of the
audience's vocabulary. That rejection is recorded with a reason; it is not
silently discarded. The author can accept F2 and F5, then revise the paragraph
in the existing editor. Humanvoice remains beside the work rather than taking
ownership of it.

\section{The author's revision}

One possible revision is:

\begin{quote}\small
"The constrained reconstruction compares the same likelihood components,
observations, and counterfactual paths as the baseline model. In the
illustrative calibration, the funding wedge is 35 basis points (Table~3), and
the comparison tests whether the solver preserves that wedge and the stated
likelihood terms (Equation~(12)). If those objects agree, the note supports a
larger comparative run; it does not by itself establish a general policy
effect."
\end{quote}

The revision makes the actors and objects visible, turns the number into a
quantity with a unit, and gives the reader a decision consequence. It also
retains a qualification because a reader could otherwise mistake a calibration
check for a policy result. That qualification now follows the positive result
instead of replacing it.

\begin{repairpair}
\before{"We do not claim that this proves the policy works in general; rather, it is a robust and transparent framework for checking the relevant objects."}
\after{"If those objects agree, the note supports a larger comparative run; it does not by itself establish a general policy effect."}
\end{repairpair}

The second sentence is shorter, but brevity is not the product's objective. The
important change is that it states what the evidence can support. The author
could choose a different wording and still answer the same audience and
decision.

\section{The protected comparison}

After the author saves the revision, the comparison engine checks the objects a
reader can least afford to lose while the prose is being polished.

The system does not freeze those objects. As Figure
\ref{fig:protected-diff} shows, it classifies their correspondence and requires
a human disposition when a consequential difference appears.

\hvFigureProtectedDiff

\begin{table}[H]
\centering
\small
\begin{tabular}{@{}p{0.31\textwidth}p{0.22\textwidth}p{0.34\textwidth}@{}}
\toprule
Protected object & Result & What the author sees \\
\midrule
Equation body & unchanged & The normalized mathematical expression is unchanged;
the surrounding explanation changed. \\
Numerical result & unchanged & 35 basis points remains 35 basis points with the
same sign and table anchor. \\
Citation and cross-reference & unchanged & Table~3 and Equation~(12) still
resolve to the same objects. \\
Claim scope & changed explicitly & "Supports a larger comparative run" is a
new, narrower conclusion and requires author confirmation. \\
Private label & removed & The source map identifies the deletion and links it to
F1. \\
\bottomrule
\end{tabular}
\caption{The diff distinguishes a safe explanatory repair from a substantive
change that needs a human decision.}
\label{tab:protected-diff}
\end{table}

This is more informative than a red-and-green visual diff alone. A visual diff
shows where characters moved; a protected comparison says which meaning-bearing
objects survived and which changed. It also creates a useful evaluation record:
the author accepted F1, F2, and F5, edited F3, and confirmed the preserved
number and equation.

\section{The reader's turn}

The final step is not another model score. A reader who has not seen the
project history receives the rendered note and answers five questions:

\begin{enumerate}
\item What comparison is being made?
\item What does the 35-basis-point quantity represent?
\item Which objects must agree before the team begins the larger run?
\item What does the paragraph not establish?
\item What action should the reader take if the checks pass?
\end{enumerate}

Suppose the reader answers the first four correctly but says that the requested
action is still unclear. That is a useful failure. The author has learned that
the paragraph is more precise but the document's decision boundary is not yet
visible. Humanvoice records a reader revision request rather than converting
the four correct answers into a passing score.

\begin{takeawaybox}{The result of one review}
For one document and one version, Humanvoice produces a source snapshot, a
small set of located findings, an author decision for each finding, a protected
before-and-after comparison, and a reader decision with elapsed time and
requested changes. The packet is inspectable by an author, measurable by an
evaluator, and priceable by a sponsor.
\end{takeawaybox}

\section{The six capabilities}

For this case, an implementer needs six connected capabilities. Each one
changes a real decision in the writing path; none is a list of packages to
install.

\begin{description}
\item[Humanization brief compiler.] Turn the reader, genre, prior knowledge,
  evidence boundary, protected-object policy, voice constraints, and named
  exemplars into a versioned brief bound to one immutable source. Critical
  blanks remain visible and block rewriting.
\item[Concept inventory and baseline.] Partition every reader-facing source
  span, record the substantive concepts it carries, assign each concept the
  explanation it owes this reader, order prerequisites, and freeze the reviewed
  result. An independent reverse pass reads the inventory back against the
  source; a model may not certify its own extraction.
\item[Source-grounded writer.] Rewrite one source passage at a time from the
  exact span, its assigned concepts and obligations, and its protected objects.
  A replaceable model may write, but it cannot drop a concept, widen a claim,
  add a citation, or skip a prerequisite without a recorded decision.
\item[Semantic preflight.] Verify the candidate revision, not just the source:
  concept correspondence over the frozen baseline, explanation fulfillment with
  located evidence, dependency order, unsupported additions, scaffolding
  decisions, and exact-object integrity. Unknown or conflicting results fail
  closed.
\item[Located repair and comparison.] Map a failed check to one causal repair,
  rebuild the affected unit, detect oscillation, and compare protected
  equations, numbers, citations, tables, quotations, and claims.
\item[Patch assembly and release.] Apply accepted replacements to a copy of the
  source tree by stable offsets, leave untouched bytes byte-identical, compile
  the revision and a whole-tree comparison, and produce a clean reader packet
  only after the gates pass. The named human still owns acceptance.
\end{description}

The system does not need to train a new text-generation model. It needs a stable
contract around a replaceable writer: the model receives only the source span it
is rewriting, the concepts it must preserve, and the evidence it may use; the
preflight controller can reject its output and request a cause-specific repair.
That distinction is what keeps fluent prose from silently losing content.

\section{The minimum useful interface}

An author starts from a finished manuscript and asks for a humanized revision:

\begin{quote}\small
\texttt{hv humanize manuscript/ --brief brief.yaml}
\end{quote}

That single command orchestrates the resumable phases, each of which can also be
run directly for inspection or recovery:

\begin{quote}\small
\texttt{hv init} \quad \texttt{hv inventory} \quad \texttt{hv plan} \quad
\texttt{hv rewrite}\par
\texttt{hv preflight} \quad \texttt{hv repair} \quad \texttt{hv assemble} \quad
\texttt{hv release}
\end{quote}

These verbs are stages of one path, not independent products. \texttt{hv init}
snapshots the immutable source and refuses a brief with unresolved critical
fields. \texttt{hv inventory} builds the concept baseline and refuses to proceed
while a reader-facing span is uncovered or a classification is unresolved.
\texttt{hv plan} assigns concepts, obligations, and rewrite units in dependency
order. \texttt{hv rewrite} rewrites one source passage and publishes nothing on
truncation or incomplete correspondence. \texttt{hv preflight} verifies the
assembled candidate. \texttt{hv repair} applies one named causal repair,
rebuilds, and stops on oscillation. \texttt{hv assemble} patches a copy of the
source tree and builds the revision and its comparison. \texttt{hv release}
creates the reader packet only when every required result is a pass. A missing
result is a failure, not a pass, and there is no generic release exception.
Greenfield authoring from a brief remains available only under an explicit
\texttt{hv legacy} namespace and cannot satisfy a release gate.

The older capture, inspect, and compare services remain useful internally, but
they no longer define the author's journey. The important property is that every
service refers to an immutable source snapshot, a frozen baseline, and a render
hash rather than to an unowned draft.

The same record should be usable from a command line, a notebook, or an editor
extension. A sponsor should not need to know which parser or linter produced a
finding in order to understand the decision. Conversely, an engineer should
be able to replace a parser without changing the reader rubric or the primary
evaluation outcome.

\section{The boundary of the first version}

The system focuses on humanizing finished manuscripts for reader understanding,
not on judging whether an individual document is ``human.''
Detector studies show high
false-positive risk for formal and non-native writing and easy evasion through paraphrase
\citep{liang2023detectors,weberwulff2023testing,krishna2023paraphrasing}. A
population-level monitoring question may be studied later, but it is outside
the revision problem described here.

The system also leaves the source untouched and in the author's hands, keeps
acceptance with the named reader, and treats style rules as observations rather
than a universal house style. Authorized humanization of a finished manuscript is
in scope; summarization, shortening to a length target, and autonomous acceptance
are not. Those choices keep the measurable outcome aligned with the reader's
decision while ensuring that unreadable prose is repaired before it consumes that
reader's time.
